\textbf{结论名称：} 和角公式（两角和与差的正弦、余弦、正切）

\textbf{适用条件：} $\alpha,\beta$ 为任意角；对 $\tan(\alpha\pm\beta)$ 需保证 $\alpha,\beta,\alpha\pm\beta \neq \frac{\pi}{2}+k\pi\,(k\in\mathbb{Z})$，且分母 $1\mp\tan\alpha\tan\beta\neq0$。

\textbf{变量说明：} $\alpha,\beta$ 为任意实数角，相关三角函数应在定义域内。

\textbf{核心公式：}
\[
\boxed{
\begin{aligned}
\sin(\alpha\pm\beta) &= \sin\alpha\cos\beta \pm \cos\alpha\sin\beta \\[2pt]
\cos(\alpha\pm\beta) &= \cos\alpha\cos\beta \mp \sin\alpha\sin\beta \\[2pt]
\tan(\alpha\pm\beta) &= \dfrac{\tan\alpha \pm \tan\beta}{1 \mp \tan\alpha\tan\beta}
\end{aligned}
}
\]

\textbf{使用场景：} 化简三角表达式、计算非特殊角三角函数值、三角恒等证明。

\textbf{简单例子：} 求 $\sin 75^\circ$。
\[
\sin 75^\circ = \sin(45^\circ + 30^\circ) = \sin45^\circ\cos30^\circ + \cos45^\circ\sin30^\circ
= \frac{\sqrt{2}}{2}\cdot\frac{\sqrt{3}}{2} + \frac{\sqrt{2}}{2}\cdot\frac{1}{2}
= \frac{\sqrt{6}+\sqrt{2}}{4},
\]
结果与已知特殊角值一致，可直接手算验证。